% Appendix C, input by main.tex.

\FloatBarrier
% ============================================================
\section{Recoveries of the Four Scenarios}\label{sec:scenarios}
% ============================================================

This appendix presents one recovery generated by the neuro-symbolic framework for each of the
four failure scenarios in the case study: Part slippage, Conveyor breakdown, UR5e-1 robot
breakdown and M1 machine breakdown. In each scenario, the framework synthesized a recovery outline
at~\gls{el}. Each accepted recovery event was then composed into a program using the
primitives of the resource that performs it. For each scenario, the recovery shown is a passing
outline with the fewest events. Table~\ref{tab:scenarios} summarizes the four recoveries. The
following subsections describe each failure and list its recovery events and primitive steps.

The composed programs were executed in the Gazebo simulation of the cell. Each run first
reproduced the failure, then executed the programs in order using the parameters resolved during
composition. The~\gls{cca} did not supervise these runs, and each run ended after the last recovery
event. All 29 recovery events across the four scenarios completed. The supplementary video
shows the four recoveries executed in Gazebo.

\begin{table}[!htb]
\centering\footnotesize
\setlength{\tabcolsep}{5pt}
\begin{tabularx}{\linewidth}{@{}Y c c c c c@{}}
\toprule
Scenario & \shortstack{Outline\\turns} & Events & \shortstack{Composition\\calls} & \shortstack{Primitive\\steps} & \shortstack{Completed\\in Gazebo} \\
\midrule
Part slippage & 7 & 4 & 9 & 16 & 4 of 4 \\
Conveyor breakdown & 12 & 7 & 15 & 23 & 7 of 7 \\
UR5e-1 robot breakdown & 11 & 4 & 19 & 20 & 4 of 4 \\
M1 machine breakdown & 20 & 14 & 49 & 62 & 14 of 14 \\
\bottomrule
\end{tabularx}
\caption{Overview of the four recoveries and their execution results in Gazebo.}\label{tab:scenarios}
\end{table}

\FloatBarrier
\subsection{Part slippage}\label{sec:part-slippage}

\code{ur5e-4} drops the small gear while placing it on the assembly board. \code{ur5e-3},
which is holding the square peg, first places that peg in the buffer. \code{ur5e-4} then moves
out of the board area, and \code{ur5e-3} picks up the gear and places it in its slot.
Appendices~A and~B describe this recovery in detail.

{\footnotesize\setlength{\tabcolsep}{4pt}
\begin{xltabular}{\linewidth}{@{}r >{\raggedright\arraybackslash}p{5.9cm} l Y@{}}
\toprule
\# & Event, as named by the~\gls{llm} & Resource & Primitive steps, in order \\
\midrule
\endfirsthead
\toprule
\# & Event, as named by the~\gls{llm} & Resource & Primitive steps, in order (continued) \\
\midrule
\endhead
\bottomrule
\endfoot
\bottomrule
\caption{Part slippage: the recovery events and their primitive steps. All 4 events completed in Gazebo.}\label{tab:part-slippage}\\
\endlastfoot
1 & \recordname{stage\_KET4\_in\_buffer} & \code{ur5e-3} & \code{compute_place_targets}, \code{move_cartesian}\,$\times$2, \code{release_part}, \code{move_cartesian} \\
2 & \recordname{retract\_ur5e4\_home\_clear\_board} & \code{ur5e-4} & \code{move_to_named_pose} \\
3 & \recordname{pick\_dropped\_gear\_small\_ur5e3} & \code{ur5e-3} & \code{compute_pick_targets}, \code{move_cartesian}\,$\times$2, \code{grasp_part}, \code{move_cartesian} \\
4 & \recordname{place\_gear\_small\_on\_board\_ur5e3} & \code{ur5e-3} & \code{compute_place_targets}, \code{move_cartesian}\,$\times$2, \code{release_part}, \code{move_cartesian} \\
\end{xltabular}}

\FloatBarrier
\subsection{Conveyor breakdown}\label{sec:conveyor}

While \code{ur5e-1} is taking the square peg to the \code{Conveyor}, the \code{Conveyor}
reports that it is unavailable. The recovery goal is to place the peg in zone~1 of the buffer,
where the \code{Conveyor} would normally carry it. \code{ur5e-1} places the peg on the M1 staging tray, and
the \code{KMR} carries it to the buffer.

{\footnotesize\setlength{\tabcolsep}{4pt}
\begin{xltabular}{\linewidth}{@{}r >{\raggedright\arraybackslash}p{5.9cm} l Y@{}}
\toprule
\# & Event, as named by the~\gls{llm} & Resource & Primitive steps, in order \\
\midrule
\endfirsthead
\toprule
\# & Event, as named by the~\gls{llm} & Resource & Primitive steps, in order (continued) \\
\midrule
\endhead
\bottomrule
\endfoot
\bottomrule
\caption{Conveyor breakdown: the recovery events and their primitive steps. All 7 events completed in Gazebo.}\label{tab:conveyor}\\
\endlastfoot
1 & \recordname{place\_on\_M1\_staging\_tray} & \code{ur5e-1} & \code{compute_place_targets}, \code{move_cartesian}\,$\times$2, \code{release_part}, \code{move_cartesian} \\
2 & \recordname{navigate\_KMR\_to\_M1} & \code{KMR} & \code{move_base} \\
3 & \recordname{ur5e1\_retract\_home} & \code{ur5e-1} & \code{move_to_named_pose} \\
4 & \recordname{kmr\_pick\_from\_M1\_staging\_tray} & \code{KMR} & \code{detect_parts}, \code{compute_pick_targets}, \code{move_to_named_pose}, \code{move_cartesian}\,$\times$2, \code{grasp_part}, \code{move_cartesian} \\
5 & \recordname{kmr\_move\_base\_to\_buffer} & \code{KMR} & \code{move_base} \\
6 & \recordname{kmr\_place\_on\_buffer\_zone\_1} & \code{KMR} & \code{compute_place_targets}, \code{move_cartesian}\,$\times$2, \code{release_part}, \code{move_cartesian}\,$\times$2 \\
7 & \recordname{kmr\_stow\_arm\_and\_return\_to\_storage} & \code{KMR} & \code{move_to_named_pose}, \code{move_base} \\
\end{xltabular}}

\FloatBarrier
\subsection{UR5e-1 robot breakdown}\label{sec:ur5e1}

The gripper of \code{ur5e-1} fails during an attempt to pick up the square peg at \code{M1}.
The recovery goal is to place the peg on the \code{Conveyor}. \code{ur5e-1} moves home, and the
\code{KMR} collects the peg at \code{M1}, delivers it to the \code{Conveyor} and returns to \code{Storage}.

{\footnotesize\setlength{\tabcolsep}{4pt}
\begin{xltabular}{\linewidth}{@{}r >{\raggedright\arraybackslash}p{5.9cm} l Y@{}}
\toprule
\# & Event, as named by the~\gls{llm} & Resource & Primitive steps, in order \\
\midrule
\endfirsthead
\toprule
\# & Event, as named by the~\gls{llm} & Resource & Primitive steps, in order (continued) \\
\midrule
\endhead
\bottomrule
\endfoot
\bottomrule
\caption{UR5e-1 robot breakdown: the recovery events and their primitive steps. All 4 events completed in Gazebo.}\label{tab:ur5e1}\\
\endlastfoot
1 & \recordname{Retract the empty arm to its named home pose} & \code{ur5e-1} & \code{move_to_named_pose} \\
2 & \recordname{Retrieve the completed part from M1 with KMR} & \code{KMR} & \code{move_to_named_pose}, \code{move_base}, \code{detect_parts}, \code{compute_pick_targets}, \code{move_cartesian}\,$\times$2, \code{grasp_part}, \code{move_cartesian}, \code{move_to_named_pose} \\
3 & \recordname{Deliver to conveyor loading position 1 using a south-facing base heading} & \code{KMR} & \code{move_to_named_pose}, \code{move_base}, \code{compute_place_targets}, \code{move_cartesian}\,$\times$2, \code{release_part}, \code{move_cartesian}\,$\times$2 \\
4 & \recordname{Return KMR to Storage and verify idle readiness} & \code{KMR} & \code{move_to_named_pose}, \code{move_base} \\
\end{xltabular}}

\FloatBarrier
\subsection{M1 machine breakdown}\label{sec:machining}

\code{M1} fails during the square trim of the square peg. The recovery goal is to complete that
trim. \code{M2} can perform the trim, but \code{M2} is already trimming the round peg.
\code{M2} pauses that job, and \code{ur5e-2} moves the round peg to M2's staging tray. The \code{KMR}
brings the square peg from \code{M1} to \code{M2}, which performs the trim. \code{ur5e-2} then places the finished
square peg on the \code{Conveyor}. Finally, \code{ur5e-2} reloads the round peg, and \code{M2} resumes its job.

{\footnotesize\setlength{\tabcolsep}{4pt}
\begin{xltabular}{\linewidth}{@{}r >{\raggedright\arraybackslash}p{5.9cm} l Y@{}}
\toprule
\# & Event, as named by the~\gls{llm} & Resource & Primitive steps, in order \\
\midrule
\endfirsthead
\toprule
\# & Event, as named by the~\gls{llm} & Resource & Primitive steps, in order (continued) \\
\midrule
\endhead
\bottomrule
\endfoot
\bottomrule
\caption{M1 machine breakdown: the recovery events and their primitive steps. All 14 events completed in Gazebo.}\label{tab:machining}\\
\endlastfoot
1 & \recordname{Pause M2\textquotesingle{}s round-part job for temporary unloading} & \code{M2} & \code{pause_job} \\
2 & \recordname{Unload the paused round part into M2\textquotesingle{}s staging tray} & \code{ur5e-2} & \code{detect_parts}, \code{compute_pick_targets}, \code{move_cartesian}\,$\times$2, \code{grasp_part}, \code{move_cartesian}, \code{compute_place_targets}, \code{move_cartesian}\,$\times$2, \code{release_part}, \code{move_cartesian} \\
3 & \recordname{Navigate to M1 and retrieve the interrupted square part} & \code{KMR} & \code{move_to_named_pose}, \code{move_base}, \code{compute_pick_targets}, \code{move_cartesian}\,$\times$2, \code{grasp_part}, \code{move_cartesian} \\
4 & \recordname{Select M2\textquotesingle{}s square program before loading} & \code{M2} & \code{select_program} \\
5 & \recordname{Return ur5e-2 home to clear M2 access} & \code{ur5e-2} & \code{move_to_named_pose} \\
6 & \recordname{Navigate to M2 and load the interrupted square part} & \code{KMR} & \code{move_to_named_pose}, \code{move_base}, \code{compute_place_targets}, \code{move_cartesian}\,$\times$2, \code{release_part}, \code{move_cartesian}\,$\times$2 \\
7 & \recordname{Navigate KMR to Storage to clear M2} & \code{KMR} & \code{move_base} \\
8 & \recordname{Start the square-part job on M2} & \code{M2} & \code{start_job} \\
9 & \recordname{Await and verify completion of square trimming on M2} & \code{M2} & \code{complete_job} \\
10 & \recordname{Transfer the completed square part from M2 to Conveyor} & \code{ur5e-2} & \code{detect_parts}, \code{compute_pick_targets}, \code{move_cartesian}\,$\times$4, \code{grasp_part}, \code{move_cartesian}, \code{compute_place_targets}, \code{move_cartesian}\,$\times$4, \code{release_part}, \code{move_cartesian} \\
11 & \recordname{Reload the paused round part from M2\textquotesingle{}s staging tray} & \code{ur5e-2} & \code{compute_pick_targets}, \code{move_cartesian}\,$\times$4, \code{grasp_part}, \code{move_cartesian}, \code{compute_place_targets}, \code{move_cartesian}\,$\times$2, \code{release_part}, \code{move_cartesian} \\
12 & \recordname{Return ur5e-2 home to clear M2 for resumed processing} & \code{ur5e-2} & \code{move_to_named_pose} \\
13 & \recordname{Restore M2\textquotesingle{}s circle program before a separate resume event} & \code{M2} & \code{select_program} \\
14 & \recordname{Resume the reloaded paused round-part job} & \code{M2} & \code{resume_job} \\
\end{xltabular}}
